\section{Software Architecture}\label{sec:architecture}

\subsection{Overview}

\diagram{architecture.png}{0.95}{System architecture: React front end, FastAPI service,
compiler core and persistence.}

The system is layered so that the compiler core has no knowledge of the web application
around it. \texttt{core/} depends on nothing but the standard library and pydantic; it can
be exercised entirely from the command line. The API layer is a thin translation between
HTTP and core function calls, and the database stores corpus records only --- never derived
analysis, which is always recomputed so it cannot go stale relative to the grammar.

\begin{center}
\begin{tabular}{@{}lll@{}}
\toprule
\textbf{Layer} & \textbf{Module} & \textbf{Responsibility} \\
\midrule
Compiler core & \texttt{core/text.py}     & scanning, canonicalisation, folding \\
              & \texttt{core/lexer.py}    & tokenisation, lexicon resolution \\
              & \texttt{core/lexicon.py}  & entry storage, lookup indices, homographs \\
              & \texttt{core/grammar.py}  & grammar model, transformations, ledger \\
              & \texttt{core/ll1.py}      & FIRST/FOLLOW, table construction \\
              & \texttt{core/parser.py}   & predictive parser \\
              & \texttt{core/suggest.py}  & repair suggestions \\
              & \texttt{core/analysis.py} & snapshots, frequencies, topics \\
              & \texttt{core/corpus.py}   & validation, revisions, provenance \\
\midrule
Service       & \texttt{api/main.py}      & application factory, middleware \\
              & \texttt{api/routes/}      & analysis, auth, corpus, export, grammar, statistics \\
\midrule
Interface     & \texttt{cli.py}           & offline analysis, seeding, export \\
              & React front end           & analyser, grammar explorer, corpus manager \\
\bottomrule
\end{tabular}
\end{center}

\subsection{Use cases}

\diagram{use-cases.png}{0.88}{Actors and use cases supported by the application.}

\subsection{Front end}

The web client is React with TypeScript and Tailwind. Three views matter for the
demonstration:

\begin{description}[leftmargin=1.4em, style=nextline]
  \item[Analyser] Type or paste a statement and see the token stream, the derivation, the
    verdict, the frequency tables and the detected topics. Individual tokens expand to show
    the canonical form, terminal class, donor language and gloss, so the lexicon is
    inspectable without leaving the analysis.
  \item[Grammar explorer] The original grammar, the transformation ledger, the transformed
    grammar, FIRST/FOLLOW sets, the LL(1) table and the full lexicon --- with incremental
    search that filters the lexicon as you type, and an explanation of what each section
    means and why the transformation was needed.
  \item[Corpus manager] Add, edit and review statements. Statement identifiers and
    collector identity are assigned server-side from the signed-in account; edits create
    revisions rather than overwriting.
\end{description}

\subsection{Deployment}

The service is containerised and runs behind a reverse proxy on a self-hosted VPS. The
image ships the built Python wheel together with the Alembic configuration and migration
scripts; development-only material such as seeding scripts is deliberately excluded, and
seeding is instead exposed as a CLI subcommand of the installed package so that it works
identically inside and outside the container.

\begin{limitbox}[Before going live]
The self-service signup code is still set to a development placeholder. It must be replaced
with a strong value in the deployment environment before the instance is exposed publicly.
\end{limitbox}
